\section{学历中下层：被系统和标签漏掉的 ADHD}

如果说学历上层是「被聪明掩盖」，那么学历中下层就是\term{被系统漏掉}——而且在被漏掉的过程中，还被打上了一层最毒的标签。

\subsection{画像}

也许从未确诊，甚至根本不知道 ADHD 是什么；早早离开学校，或勉强读完技校、大专；做着需要动手、需要即时反馈、不需要长时间线性规划的工作；也可能长期在家，状态低迷，被家人说「懒」、被自己说「废」。

他们中的很多人，在互联网上第一次看到 ADHD 的描述时，是哭着看完的——不是因为确诊了，而是因为终于有一个词，能解释自己这么多年的「不对劲」。

\subsection{他们经历了什么}

理解学历中下层的 ADHD，必须从学校说起。因为\term{学校可能是全世界对 ADHD 最不友好的环境}：

\begin{itemize}
  \item \term{教育系统默认深度优先算法}：一节课四十五分钟、一章接一章、知识线性积累、一次没跟上就永远跟不上。这恰好是 ADHD 算法最贵的运行模式，和它是天生的仇家。
  \item \term{早期失败，归因到态度而不是大脑}：「不认真」「懒」「笨」「态度有问题」。归因到态度，比归因到大脑更伤人——因为它意味着「你本来可以，是你不愿意」，是把「不能」翻译成了「不愿」，再翻译成「人品问题」。
  \item \term{自我概念崩坏}：在最重要的成长期，被反复告知「你不行」，于是「我做不到」成了默认信念，哪怕他其实很聪明。
  \item \term{回避 → 更多失败}：逃避所有可能暴露「我不行」的场合——不交作业、不去考试、不竞选、不发言——于是错过更多机会，失败进一步坐实。
  \item \term{家庭资源不足}：没有诊断的机会、没有好老师、没有「再试一次」的试错资本。学历的差距会滚雪球：一步落下，后面每一步都在为前面买单。
\end{itemize}

\subsection{习惯观察}

学历中下层的 ADHD，在外人眼里常常是这样：

\begin{checklist}
  \item \textbf{三分钟热度}：兴趣一来就上头，几天几夜不睡觉；兴趣一退就弃坑，换下一个。频繁换工作、频繁开始又放弃；
  \item \textbf{用即时刺激填满生活}：游戏、短视频、烟酒——在旁人看来是堕落，其实是一个燃料不足的大脑在拼命找刺激自救；
  \item \textbf{隐藏技能}：极强的动手能力、直觉、口才、销售天赋、危机处理能力——这些在学校里不被计分，在社会上也不被当成「学历」，但它们真实存在；
  \item \textbf{「地下」的成功模式}：很多人最终在不需要学历的赛道里反而活得不错——手艺、销售、自由职业、创业。这正是广度算法的用武之地。
\end{checklist}

\subsection{成因：不是能力不够，是失败的积累压垮了执行功能}

要特别强调一件事：\term{学历中下层的 ADHD，核心问题从来不是「能力」}。同一颗大脑，放进一个有人推着走的环境，能读出来；放进一个没人管、没人信的环境，可能连毕业都难。

他们真正缺的是两样：\term{一次成功的经验}（「我能做到」的证据），和\term{一个不评判他的环境}。前者是燃料，后者是脚手架。这两样都没有，执行功能就会在「反正我也做不到」的预测里彻底躺平。

\begin{insight}
学历中下层最缺的不是方法，是\term{「我能做到」的证据}。方法可以学，但若心里那个「我不行」的信念不破，什么方法都会在第一次失败时被当成「你看，果然不行」。所以他们的自救顺序，和学历上层正好相反：学历上层先补系统，学历中下层先修自我概念。
\end{insight}

\subsection{一句话总结这一章}

学历中下层的核心矛盾：\term{不是做不到，是相信自己做不到}。他们最需要的是被给一次「第二次机会」——不管是自己给的，还是世界给的。

而这一章写在这里，就是那份第二次机会的第一步：\term{你的学历低，不代表你不行；它只代表，你一直没有脚手架。}
